\usetikzlibrary{shapes, backgrounds}

\chapter{Matroidi}

\section{Definizione}
Un Sistema di Indipendenza $\langle E,F \rangle$ è un matroide se e solo se vale la seguente proprietà (detto "di scambio")

{
    \large
    $$
    \forall A,B \in F \text{ con } |B| > |A| \quad \exists b \in B \backslash A \text{ tale che } A \cup \{b\} \in F
    $$

    $$
    |B| > |A| \text{ si può sostituire equivalentemente con } |B| = |A|+1 
    $$
}

\subsection{Esempi di matroidi}
\begin{itemize}
    \item Il S.I $\langle E, \mathcal{P}(E) \rangle$ dell'esempio 1 è un matroide
    \item Il S.I dell'esempio 2 è un matroide
    \item Il S.I dell'esempio 3 (matroide grafico) è un matroide
    \item Il S.I dell'esempio 4 non è un matroide
    $$
    \begin{aligned}
        A = \{x_1\} \in F & \quad (\text{ing}(A) = 30 \leq 45) \\
        B = \{x_2, x_3\} \in F & \quad (\text{ing(B)} = 45 \leq 45) \\
        \not \exists b \in B \backslash A : A \cup \{b\} \in F & \quad \text{Non rispetta la proprietà di scambio } \\ 
    \end{aligned}
    $$
    \begin{center}
        Infatti $ing(A \cup b) > 45$ con $b \in B \backslash A$
    \end{center}
\end{itemize}

\section{Teorema della dimensione di un matroide}

\subsection{Proposizione}
Sia $\langle E,F \rangle$ un matroide, tutti gli insiemi massimali di $\langle E,F \rangle$ hanno la stessa cardinalità

\subsection{Dimostrazione}
Per assurdo, siano $A,B \in F$ massimali con $|B| > |A|$ per la proprietà dello scambio 

$$
\exists b \in B \backslash A : A \cup \{b\} \in F
$$

quindi $A$ è estendibile, il che è una contraddizione

\newpage
\section{Matroidi grafici}

\subsection{Obiettivo}
Vogliamo ora far vedere che il matroide grafico è un matroide.

\begin{itemize}
    \item Sia $G = (V,E)$ un grafo orientato
    \item Sia $\langle E,F \rangle$ il S.I con $F=\{A \subseteq E : (V,A) \text{ è una foresta}\}$
\end{itemize}

\begin{figure}[htbp]
\centering
\definecolor{highlightpink}{RGB}{255, 204, 218}
\definecolor{softpurple}{RGB}{220, 210, 245}
\definecolor{softblue}{RGB}{210, 220, 250}

\begin{tikzpicture}[
    vertex/.style={circle, draw, thick, inner sep=0pt, minimum size=6.5mm, font=\small\bfseries},
    hl/.style={draw=highlightpink, line width=7pt, line cap=round, line join=round},
    edge/.style={thick, line cap=round},
    textnode/.style={anchor=west, font=\normalsize}
]

% ==================== INTESTAZIONE ====================
\node[anchor=west] at (0, 0.8) {\textbf{ESEMPIO:}\quad $G = (V, E)$\quad $V = \{1, 2, 3, 4, 5, 6\}$};

% ==================== RIGA 1 (Grafo base G) ====================
\begin{scope}[yshift=0cm]
    % Nodi
    \node[vertex] (1) at (0, 0) {1};
    \node[vertex] (2) at (1.8, 0) {2};
    \node[vertex] (3) at (3.6, 0) {3};
    \node[vertex] (4) at (0, -1.3) {4};
    \node[vertex] (5) at (1.8, -1.3) {5};
    \node[vertex] (6) at (3.6, -1.3) {6};

    % Archi
    \draw[edge] (1) -- (2);
    \draw[edge] (2) -- (3);
    \draw[edge] (1) -- (4);
    \draw[edge] (4) -- (2);
    \draw[edge] (4) -- (5);
    \draw[edge] (5) -- (6);

    % Testo
    \node[textnode] at (4.5, -0.65) {$E = \{(1,2), (1,4), (2,3), (2,4), \dots\}$};
\end{scope}

% ==================== RIGA 2 (Foresta A) ====================
\begin{scope}[yshift=-2.8cm]
    % Evidenziazione archi
    \begin{pgfonlayer}{background}
        \draw[hl] (0, 0) -- (0, -1.3);
        \draw[hl] (0, -1.3) -- (1.8, 0);
        \draw[hl] (0, -1.3) -- (1.8, -1.3);
    \end{pgfonlayer}

    % Nodi
    \node[vertex] (1) at (0, 0) {1};
    \node[vertex] (2) at (1.8, 0) {2};
    \node[vertex, fill=softpurple] (3) at (3.6, 0) {3};
    \node[vertex] (4) at (0, -1.3) {4};
    \node[vertex] (5) at (1.8, -1.3) {5};
    \node[vertex, fill=softblue] (6) at (3.6, -1.3) {6};

    % Archi
    \draw[edge] (1) -- (4);
    \draw[edge] (4) -- (2);
    \draw[edge] (4) -- (5);

    % Testo
    \node[textnode] at (4.5, 0) {$A = \{(1,4), (2,4), (4,5)\}$};
    \node[textnode] at (4.5, -0.65) {$(V, A)$ \`e una foresta (di 3 alberi)};
    \node[textnode] at (4.5, -1.3) {$A \in \mathcal{F}$};
\end{scope}

% ==================== RIGA 3 (Foresta B) ====================
\begin{scope}[yshift=-5.6cm]
    % Evidenziazione archi
    \begin{pgfonlayer}{background}
        \draw[hl] (0, 0) -- (0, -1.3);
        \draw[hl] (0, 0) -- (1.8, 0);
        \draw[hl] (1.8, 0) -- (3.6, 0);
        \draw[hl] (1.8, -1.3) -- (3.6, -1.3);
    \end{pgfonlayer}

    % Nodi
    \node[vertex] (1) at (0, 0) {1};
    \node[vertex] (2) at (1.8, 0) {2};
    \node[vertex] (3) at (3.6, 0) {3};
    \node[vertex] (4) at (0, -1.3) {4};
    \node[vertex] (5) at (1.8, -1.3) {5};
    \node[vertex] (6) at (3.6, -1.3) {6};

    % Archi
    \draw[edge] (1) -- (4);
    \draw[edge] (1) -- (2);
    \draw[edge] (2) -- (3);
    \draw[edge] (5) -- (6);

    % Testo
    \node[textnode] at (4.5, 0) {$B = \{(1,2), (1,4), (2,3), (5,6)\}$};
    \node[textnode] at (4.5, -0.65) {$(V, B)$ \`e una foresta (di 2 alberi)};
    \node[textnode] at (4.5, -1.3) {$B \in \mathcal{F}$};
\end{scope}

% ==================== RIGA 4 (Non foresta A') ====================
\begin{scope}[yshift=-8.4cm]
    % Evidenziazione ciclo
    \begin{pgfonlayer}{background}
        \draw[hl] (0, 0) -- (1.8, 0) -- (0, -1.3) -- cycle;
    \end{pgfonlayer}

    % Nodi
    \node[vertex] (1) at (0, 0) {1};
    \node[vertex] (2) at (1.8, 0) {2};
    \node[vertex] (3) at (3.6, 0) {3};
    \node[vertex] (4) at (0, -1.3) {4};
    \node[vertex] (5) at (1.8, -1.3) {5};
    \node[vertex] (6) at (3.6, -1.3) {6};

    % Archi
    \draw[edge] (1) -- (2);
    \draw[edge] (2) -- (4);
    \draw[edge] (4) -- (1);
    \draw[edge] (5) -- (6);

    % Testo
    \node[textnode] at (4.5, 0) {$A' = \{(1,2), (1,4), (2,4), (5,6)\}$};
    \node[textnode] at (4.5, -0.65) {$(V, A')$ non \`e una foresta $\longrightarrow$ \colorbox{highlightpink}{perch\'e c'\`e un ciclo nel grafo} $(1, 2, 4)$};
    \node[textnode] at (4.5, -1.3) {$A' \notin \mathcal{F}$};
\end{scope}

\end{tikzpicture}
\end{figure}

\subsection{Lemma sul numero di componenti connesse di una foresta}

\subsubsection{Enunciato}
Sia $(V,C)$ una foresta essa contiene $t$ alberi, dove $t = |V| - |C|$

\subsubsection{Dimostrazione}
Siano $v_i / e_i$ il numero di vertici / archi dell'albero $i$ nella foresta $(V,C)$ si ha che:

{
    \large
    $$
    |C| = \sum_{i=1}^{t} e_i = \sum_{i=1}^{t} (v_i - 1) = \sum_{i=1}^{t} v_i - \sum_{i=1}^{t} 1 = |V| - t \implies t = |V| - |C|
    $$
}

\newpage
\subsection{Teorema dell'esistenza del matroide grafico}

\subsubsection{Enunciato}

\begin{itemize}
    \item Sia $(V,E)$ un grafo non orientato
    \item Sia $\langle E,F \rangle$ con $F=\{A \subseteq E : (V,A) \text{ è una foresta}\}$
\end{itemize}

Vale che $\langle E,F \rangle$ è un matroide

\subsubsection{Dimostrazione}
Occorre mostrare che:
\begin{itemize}
    \item $\forall A,B \in F$ con $|B| > A \quad \exists b \in B \backslash A : A \cup \{b\} \in F$
    \item Quindi $(V,A)$ e $(V,B)$ sono due foreste
    \begin{itemize}
        \item $(V,A)$ consiste di $|V|-|A|$ alberi $\implies$ +alberi
        \item $(V,B)$ consiste di $|V|-|B|$ alberi $\implies$ -alberi
    \end{itemize}
\end{itemize}

Allora $(V,B)$ deve contenere un albero $T$ in cui ci sono due vertici $u,v$ con $(u,v) \in B$ e che si trovano in due alberi di $A \implies A \cup \{(u,v)\} \in F$

\newpage
\section{Teorema di Rado}
\subsection{Enunciato}
Sia $\langle E,F \rangle$ un Sistema di indipendenza.

Vale che $\langle E,F \rangle$ è un matroide se e solo se $\forall w : E \to R^+$ Greedy-Max($E,F,w$) risolve il problema di massimo associato a $\langle E,F \rangle$ e $w$, ossia ritorna un massimale di peso massimo tra tutti i massimali di $F$

\subsection{Dimostrazione}
Supponiamo che $\langle E, F \rangle$ è un matroide.

\begin{itemize}
    \item Sia $w:E \to R^+$ una qualsiasi funzione peso 
    \item Sia $S = \{b_1, \dots, b_p\}$ con $w(b_1) \geq \dots w(b_p)$ l'insieme ritornato da Greedy-Max($E,F,w$)
\end{itemize}

Sappiamo che $S \in F$ e $S$ è massimale.

Sia ora $A \in F$ con $A$ massimale \uline{dobbiamo far vedere che $w(S) \geq w(A)$}

Siccome $\langle E,F \rangle$ è un matroide allora $|A|=|S|=p$ (tutti i massimali hanno la stessa cardinalità per il teorema della dimensione di un matroide)

Quindi possiamo scrivere $A = \{a_1, \dots, a_p\}$ assumendo che $w(a_1) \geq \dots \geq w(a_p)$ senza perdità di generalità (WLONG).

Mostriamo ora che \uline{$\forall i \in \{1, \dots, p\}\ w(b_i) \geq w(a_i)$} fatto che consente di concludere che $\mathbf{w(S) \geq w(A)}$.

Se per assurdo \uline{$\exists i \in \{1, \dots, p\}$ tale che $w(b_i) \leq w(a_i)$ sia $k$ il più piccolo di tali $i$.}

Quindi $\mathbf{w(b_k) < w(a_k)}$

Consideriamo gli insiemi
{
    \large
    
    \begin{align*}
        & D = \{b_1, \dots, b_{k-1}\} \subseteq S \to D \in F \\
        & D' = \{a_1, \dots, a_{k-1}, a_k\} \subseteq A \to D' \in F 
    \end{align*}
}

Siccome $|D'| > |D|$ per la proprietà dello scambio

{
    \large
    $$
        \exists a_j \in D' \backslash D : D \cup \{a_j\} \in F
    $$
}

Siccome $D$ è costituito dai primi $k-1$ elementi di $S$, quindi i primi $k-1$ elemeti selezionati da Greedy-Max e il $k$-esimo selezionato è $b_k$, essendo \uline{$a_i \not \in D$ vale che $w(b_k) \geq w(s_i)$}.

Inoltre, \uline{siccome $a_i \in D'$, vale che $w(a_i) \geq w(a_k)$} quindi da $w(b_k) \geq w(s_i)$ e $w(a_i) \geq w(a_k)$ \textbf{che contraddice} $\mathbf{w(b_k) < w(a_k)}$

\newpage
Greedy-Max($E,F,w$) ritorna un insieme $S$ nel quale:
\begin{itemize}
    \item Vi sono tutti gli elementi di $A$ (poichè sono quelli di peso $>$ e poichè $A \in F$)
    \item Non vi è alcun elemento di $B \backslash A$ (poichè $\forall b \in B \backslash A \quad A \cup \{b\} \not \in F$)
    \item Potrebbero esserici eventualmente elementi di $(A \cup B)^C$ (questi hanno peso nullo)
\end{itemize}

$$\text{Quindi } w(S) = w(A) = |A| \cdot (1 + \frac{1}{2|A|}) = |A| + \frac{1}{2}$$

Calcoliamo ora $w(B)$, sapendo che $B = (B \backslash A) \cup (A \cap B)$

{
    \large
    \[
    \omega(B) = \underbrace{1 + \dots + 1}_{\text{un } 1 \quad \forall x \in B \setminus A} 
    + \underbrace{1 + \frac{1}{2|A|} + \dots + 1 + \frac{1}{2|A|}}_{\text{un contributo di } 1 + \frac{1}{2|A|} \quad \forall x \in A \cap B}
    \]
}

Pertanto $w(B) \geq |B| = |A| + 1 > |A| + \frac{1}{2} = w(S)$.

{
    \large
    $$
    \mathbf{\text{In definitiva } w(B) > w(S)}
    $$
}

Quindi l'insieme $S$ restituto da Greedy-Max non è di peso massimo.